\documentclass[10pt]{article}
\usepackage[margin=0.8in]{geometry}
\usepackage[T1]{fontenc}
\usepackage{lmodern}
\usepackage{hyperref}
\usepackage{booktabs}
\hypersetup{hidelinks}
\setlength{\emergencystretch}{2em}
% Generated by analyze.py from cells.tsv; do not edit.
\expandafter\def\csname Wbase\endcsname{1}
\expandafter\def\csname Nbase\endcsname{20}
\expandafter\def\csname Wprobe\endcsname{0}
\expandafter\def\csname Nprobe\endcsname{19}
\expandafter\def\csname Wsmax8\endcsname{0}
\expandafter\def\csname Nsmax8\endcsname{10}
\expandafter\def\csname Wsmax12\endcsname{0}
\expandafter\def\csname Nsmax12\endcsname{20}
\expandafter\def\csname Wsm5\endcsname{0}
\expandafter\def\csname Nsm5\endcsname{10}
\def\Wcommonbase{1}\def\Wcommonprobe{0}\def\Wonlybase{1}\def\Wonlyprobe{0}\def\Wcommon{19}
\expandafter\def\csname Lcrashbase\endcsname{8.550}
\expandafter\def\csname Lcrashprobe\endcsname{8.550}
\expandafter\def\csname Lcrashsmax8\endcsname{8.550}
\expandafter\def\csname Lcrashsmax12\endcsname{8.550}
\expandafter\def\csname Lcrashsm5\endcsname{10.066}
\expandafter\def\csname Lsmalldoublecrashbase\endcsname{13.983}
\expandafter\def\csname Lsmalldoublecrashprobe\endcsname{13.983}
\expandafter\def\csname Lsmalldoublecrashsmax8\endcsname{16.940}
\expandafter\def\csname Lsmalldoublecrashsmax12\endcsname{22.843}
\expandafter\def\csname Lsmalldoublecrashsm5\endcsname{7.596}
\def\Verdict{not rejected; rule 1 inconclusive}\def\RuleOne{inconclusive}\def\RuleTwo{pass}\def\RuleThree{pass}

\newcommand{\R}[1]{\csname #1\endcsname}
\title{A Direct Probe Before Suspicion Expiry:\\A Bounded Lightpool Screening Study}
\author{Claude (AI research and implementation draft)}
\date{26 November 2026 --- Unreviewed research draft}
\begin{document}
\maketitle
\begin{abstract}
A suspicion that collects no confirmations runs to a fixed maximum, and
an observer that missed the suspect's refutation then declares a reachable
member dead. Above 32 nodes the default anti-entropy interval is longer
than that maximum. We test, as an inert patch to Lightpool, one direct
probe sent one \texttt{ProbeTimeout} before the unchanged deadline, whose
acknowledgment restarts the suspicion. A predeclared campaign of 208
virtual cells compared it with the baseline and three timeout retunes;
207 completed and one was interrupted by its supervisor and retained.
On 20 held-out aged 64-node replicates the baseline made
\R{Wbase} working-path expiry death, fewer than the 3 the rules required,
so the benefit test is inconclusive. The probe and the
\texttt{SuspicionMaxMult} retunes had none. The probe left crash
detection latency unchanged to the nanosecond in both crash fixtures,
added at most 1.7\% to mean datagram counts, and passed every null
control. Every retune slowed detection by 18--63\% in one crash fixture.
The candidate is neither rejected nor admitted. This is a bounded virtual
screen, not an ordinary-network claim.
\end{abstract}

\section{Question}
In Lightpool, as in memberlist, a suspicion that collects no confirmations
runs to a fixed maximum and the observer then declares the member dead
without consulting its own path to it. An observer that missed the
member's refutation keeps the old incarnation; its later successful probes
of the member change nothing, and only the next push--pull exchange
repairs the belief. With defaults the lone maximum is 36.1\,s at 32 nodes
and 43.3\,s at 64, while the exchange interval is 30\,s up to 32 nodes and
60\,s from 33 to 64. Above 32 nodes the timer can therefore win the race
and declare a reachable member dead. We ask whether one direct probe, sent
one \texttt{ProbeTimeout} before the unchanged deadline, removes those
deaths without slowing crash detection or adding traffic, and whether a
single timeout retune does as well.

\section{Related work}
SWIM~\cite{swim} adds a suspicion subprotocol to its randomized probing
detector. Section~4.2 states that if a member ``successfully pings a
suspected member \ldots\ during the due course of the basic SWIM protocol,
it un-marks the previous suspicion'' and spreads an Alive message; entries
that time out before an Alive message is received are declared faulty.
Lightpool inherits memberlist's rule instead, under which only the
suspected member may end a suspicion, by refuting with a higher
incarnation, and Lifeguard's confirmation-scaled timeout~\cite{lifeguard}.
The candidate sits between the two: it consults the observer's own path
once, just before the deadline, and an answer only postpones the
declaration. We read the SWIM text for this draft; we have not surveyed
other implementations and make no novelty claim.

\section{Candidate and arms}
Each suspicion keeps its timer and deadline. A second timer is armed
\texttt{ProbeTimeout} before the deadline whenever the suspicion timer is
armed or re-armed, unless an expiry probe is outstanding or the deadline
is not in the future. When it fires and the member is still suspect, the
observer sends one datagram to the member's address: a ping and the held
accusation, from itself, so a target that never heard the suspicion can
refute. There is no indirect leg, TCP fallback or retry, and a failed send
is not escalated. An acknowledgment from that address restarts the
suspicion at its current confirmation level; the member stays suspect.
Silence lets the unchanged deadline declare dead. Health, probe statistics
and results are untouched. The patch changes \texttt{core/membership.go}
and \texttt{core/probe.go} by 38 lines and adds no wire message,
configuration field or journal kind.

The arms are the baseline at the pinned tooling commit (\texttt{base}),
the candidate (\texttt{probe}), \texttt{SuspicionMaxMult} 8 and 12
(\texttt{smax8}, \texttt{smax12}) and \texttt{SuspicionMult} 5
(\texttt{sm5}), the retunes applied as protocol overlays on the baseline.
The baseline defaults are \texttt{SuspicionMult} 4 and
\texttt{SuspicionMaxMult} 6.

\section{Methods}
All runs are deterministic virtual simulations with 1\,ms one-way latency,
fixed seeds derived from the replicate index, and go1.26.3 linux/amd64.
The primary workload, \texttt{correlated-dense-64-aged}, has 64 nodes; for
16.384\,s from 95\,s, datagrams from 48 nodes A to a target T are dropped
and streams fail, and the run ends at 200\,s. The 95\,s warm-up puts every
node's first periodic exchange before the fault. Its healthy control has
no faults. \texttt{crash} kills one non-bootstrap node of 32 at 15\,s;
\texttt{small-double-crash} kills two of four, where confirmations are
scarce and the suspicion timeout is near its maximum. Seven existing
32-node fixtures at development indices 1--3 serve as regression and null
controls.

The arms, 208 cells, metrics and rejection rules were committed before the
first held-out cell (indices 101 and up; indices 1--23 were used to design
the candidate). Every summary was reduced again offline by the baseline
binary and had to match. \emph{W} counts dead declarations caused by local
suspicion expiry about a running target, with neither endpoint paused and
no configured fault between them at declaration time. \emph{L} is, per
crash, the largest completed per-observer dead latency; a cell's L is the
mean over its crashes and an arm's L the mean over replicates. \emph{C} is
the attempted datagram count over the measured window. Replicates are the
unit of comparison; observer pairs within a run are dependent and are
never pooled.

The candidate is rejected if (1) baseline has at least 3 W deaths on the
20 aged replicates and the candidate more than half as many (fewer than 3
is inconclusive); (2) a retune has W no greater than the candidate's on
their common aged replicates and mean L within 5\% of the candidate's in
both crash fixtures; or (3) the candidate's mean L exceeds baseline's by
more than 5\% in either crash fixture, its mean datagram count exceeds
baseline's by more than 2\% in any fixture, it declares a running,
unpaused target dead in a cell where baseline does not, or a null control
differs.

\section{Results}
All 208 declared cells were attempted. The supervising shell's two-hour
limit killed the first full launch while recording
\texttt{probe-correlated-dense-64-aged-r117}; that attempt is retained
with its partial output and was not retried, so replicate 117 is
inconclusive for probe comparisons. The ten cells never launched were then
recorded by a continuation that refuses any already attempted cell. An
earlier launch stopped after one development cell because its evidence
directory did not exist; that attempt is retained too. Every recorded
summary equals the baseline binary's offline reduction. The verdict is
\emph{\Verdict}: rule~1 \RuleOne, rule~2 \RuleTwo, rule~3 \RuleThree.

\paragraph{Working-path deaths.}
\begin{center}
\begin{tabular}{lrrr}
\toprule
Arm & Replicates & Complete & W \\
\midrule
\texttt{base} & 101--120 & \R{Nbase} & \R{Wbase} \\
\texttt{probe} & 101--120 & \R{Nprobe} & \R{Wprobe} \\
\texttt{smax8} & 101--110 & \R{Nsmax8} & \R{Wsmax8} \\
\texttt{smax12} & 101--120 & \R{Nsmax12} & \R{Wsmax12} \\
\texttt{sm5} & 101--110 & \R{Nsm5} & \R{Wsm5} \\
\bottomrule
\end{tabular}
\end{center}
The single baseline death, in replicate 115, is absent under the probe and
\texttt{smax12}. There the probe turned 4.622\,s of dead pair-time into
suspect time and sent 6 fewer datagrams than baseline. On the
\Wcommon{} replicates where both completed, \Wonlybase{} replicate had a
death under baseline only and \Wonlyprobe{} under the probe only. The
design screen had seen 3 such deaths in 10 development replicates of the
same workload; the held-out rate was lower, and rule~1 needs at least 3.
In 19 of 20 aged replicates every arm produced the same suspect time:
each suspicion was refuted before any timeout bound mattered, so the arms
differ only where a stale lone suspicion runs out.

\paragraph{Crash detection.}
\begin{center}
\begin{tabular}{lrrrrr}
\toprule
Fixture & \texttt{base} & \texttt{probe} & \texttt{smax8} & \texttt{smax12} & \texttt{sm5} \\
\midrule
\texttt{crash}, mean L (s) & \R{Lcrashbase} & \R{Lcrashprobe} & \R{Lcrashsmax8} & \R{Lcrashsmax12} & \R{Lcrashsm5} \\
\texttt{small-double-crash}, mean L (s) & \R{Lsmalldoublecrashbase} & \R{Lsmalldoublecrashprobe} & \R{Lsmalldoublecrashsmax8} & \R{Lsmalldoublecrashsmax12} & \R{Lsmalldoublecrashsm5} \\
\bottomrule
\end{tabular}
\end{center}
The probe's per-cell maximum latencies equal baseline's in all 16 crash
replicates. The \texttt{SuspicionMaxMult} retunes match baseline at 32
nodes, where confirmations bring timers to their minimum, and cost 21\%
and 63\% when two of four nodes crash and confirmations are scarce.
\texttt{SuspicionMult} 5 is 18\% slower in \texttt{crash} and 46\% faster in
the double crash: with four nodes its confirmation target exceeds the
two possible confirmers, so every suspicion uses the minimum timeout. No retune meets rule~2's 5\% latency bound in both
fixtures, so rule~2 cannot reject the candidate whatever the death counts.

\paragraph{Cost and controls.}
Mean attempted datagram counts under the probe were within 2\% of
baseline in every fixture. The largest ratios were 1.0170 in
\texttt{threshold}, 1.0079 in \texttt{small-double-crash} and 1.0042 in
\texttt{crash}; the aged fixture was 1.0000. Traces for
\texttt{healthy}, \texttt{loss}, \texttt{correlated} and
\texttt{partial-partition} equal baseline after the header in all 12
cells. The probe produced no dead declaration about a running, unpaused
target in a cell where baseline had none. In \texttt{correlated-dense}
development index 2 it removes the known death, matching the screened
one-timer shape.

\section{Limitations}
This is a screen with few expected events, not an estimate: we report
counts and discordant replicates and no interval. The workloads are
uncalibrated virtual fixtures chosen to expose one race; they say nothing
about ordinary networks, 128 nodes or steady-state rates. The candidate
restarts the suspicion on every acknowledgment, so a member that keeps
answering this observer is never declared dead by this observer's timer
alone; that conflicts with a finite maximum timeout and is left to an
adoption decision. The simulator delivers cancelled timers by design, so
the candidate adds trace events wherever a suspicion starts even when no
expiry probe fires; one simulator golden test pinned to baseline event
counts fails under the patch for that reason.

\section{Reproduction}
The recipe \texttt{run.sh NEW\_DIR [CELL \ldots]} in
\texttt{experiments/2026-11-26-expiry-probe/} rebuilds both binaries from the pinned commit and refuses any cell whose
row differs from \texttt{cells.tsv}. \texttt{evidence-store.sh fetch
2026-11-26-expiry-probe} restores the original attempts;
\texttt{analyze.py} applies the rules to \texttt{cells.tsv};
\texttt{render.sh} renders this draft.

\begin{thebibliography}{9}
\bibitem{swim} A.~Das, I.~Gupta and A.~Motivala. SWIM: Scalable
Weakly-consistent Infection-style Process Group Membership Protocol.
\emph{DSN}, 2002. Section~4.2 read 2026-11-26 from
\url{https://www.cs.cornell.edu/projects/Quicksilver/public_pdfs/SWIM.pdf}.
\bibitem{lifeguard} A.~Dadgar, J.~Phillips and J.~Currey. Lifeguard: Local
Health Awareness for More Accurate Failure Detection. arXiv:1707.00788v2,
2018.
\end{thebibliography}
\end{document}
